\section{Efficient MC For Value-At-Risk}

\subsection{VaR And Delta-Gamma}

\subsubsection*{Loss and VaR}

Let $S$ be market prices/rates, $\Delta t$ the risk horizon, $\Delta S$ the risk-factor move, and $V(S,t)$ the portfolio value. The horizon loss is

$L=-\Delta V=V(S,t)-V(S+\Delta S,t+\Delta t),\qquad F_L(x)=P(L<x)$.

For tail probability $p$, VaR $x_p$ satisfies
$1-F_L(x_p)=P(L>x_p)=p$.

For 99\% VaR, $p=0.01$. VaR is not coherent because subadditivity can fail.

\subsubsection*{RiskMetrics}

Assume $\Delta S\sim N(0,\Sigma_S)$ and linear portfolio response $\Delta V=\delta^\top\Delta S$. Then

$L\sim N(0,\sigma_L^2),\qquad \sigma_L^2=\delta^\top\Sigma_S\delta$.

Thus 99\% VaR is approximately $2.33\sigma_L$.

\subsubsection*{Delta-gamma approximation}

For nonlinear portfolios,
$\Delta V\approx\frac{\partial V}{\partial t}\Delta t+\delta^\top\Delta S+\frac12\Delta S^\top\Gamma\Delta S$,

$\delta_i=\frac{\partial V}{\partial S_i},\qquad \Gamma_{ij}=\frac{\partial^2V}{\partial S_i\partial S_j}$.

Let $\Delta S=CZ$, $Z\sim N(0,I)$, $CC^\top=\Sigma_S$. Then

$L\approx a-(C^\top\delta)^\top Z-\frac12Z^\top(C^\top\Gamma C)Z,\qquad a=-(\Delta t)\frac{\partial V}{\partial t}$.

Choose $C$ to diagonalize the quadratic term. If $-\frac12\tilde C^\top\Gamma\tilde C=U\Lambda U^\top$, $C=\tilde C U$, $\Lambda=\operatorname{diag}(\lambda_1,\ldots,\lambda_m)$, and $b=-C^\top\delta$, then

$Q=a+b^\top Z+Z^\top\Lambda Z=a+\sum_{j=1}^m(b_jZ_j+\lambda_jZ_j^2)$.

Use $P(L>x)\approx P(Q>x)$.

\subsubsection*{Distribution of $Q$}

The cumulant generating function $\psi(\theta)=\log\mathbb{E}[e^{\theta Q}]$ is

$\psi(\theta)=a\theta+\frac12\sum_{j=1}^m\left(\frac{\theta^2b_j^2}{1-2\theta\lambda_j}-\log(1-2\theta\lambda_j)\right)$,

valid when $\max_j\theta\lambda_j<1/2$. The characteristic function is $\hat\phi(u)=e^{\psi(iu)}$. Transform inversion gives

$P(Q\le x)-P(Q\le x-y)=\frac1\pi\int_0^\infty\operatorname{Re}\left(\hat\phi(u)\left[\frac{e^{iuy}-1}{iu}\right]e^{-iux}\right)du$.

Choose $y$ large so $P(Q\le x-y)$ is negligible.

\subsubsection*{Plain MC and quantile}

Ordinary MC generates market moves, revalues the portfolio, and estimates

$P(L>x)\approx\frac1n\sum_{i=1}^n\mathbf{1}_{\{L_i>x\}}$.

The bottleneck is portfolio revaluation. With empirical CDF

$\hat F_{L,n}(x)=\frac1n\sum_{i=1}^n\mathbf{1}_{\{L_i\le x\}}$,
estimate VaR by $\hat x_p=\hat F_{L,n}^{-1}(1-p)$.

\subsection{Variance Reduction}

\subsubsection*{Control variate}

Evaluate both $L_i$ and $Q_i$ on each path. A control estimator for $P(L>x)$ is

$1-\hat F_L^{CV}(x)=\frac1n\sum_{i=1}^n\mathbf{1}_{\{L_i>x\}}-\hat\beta\left(\frac1n\sum_{i=1}^n\mathbf{1}_{\{Q_i>y\}}-P(Q>y)\right)$.

$P(Q>y)$ is known by transform inversion, and $y$ need not equal $x$. The method can disappoint in the far tail: $L$ and $Q$ may be highly correlated, but their exceedance indicators may not be, and few exceedances make $\hat\beta$ noisy.

\subsubsection*{Importance sampling}

Large $Q$ occurs when $b_jZ_j$ is large or when $\lambda_jZ_j^2$ is large with $\lambda_j>0$. Define the twisted measure
$\frac{dP_\theta}{dP}=e^{\theta Q-\psi(\theta)}$.

Then
$P(L>x)=\mathbb{E}_\theta\left[e^{-\theta Q+\psi(\theta)}\mathbf{1}_{\{L>x\}}\right]$.

Under $P_\theta$, $Z$ remains normal with diagonal covariance:

$Z\sim N(\mu(\theta),\Sigma(\theta)),\qquad \mu_j(\theta)=\frac{\theta b_j}{1-2\lambda_j\theta},\qquad \sigma_j^2(\theta)=\frac1{1-2\lambda_j\theta}$.

For $\theta>0$, the mean shifts in the sign of $b_j$ and variance increases when $\lambda_j>0$.

\subsubsection*{IS algorithm}

Choose $\theta>0$ with finite $\psi(\theta)$. For each path, generate $Z\sim N(\mu(\theta),\Sigma(\theta))$, compute $Q$, set $\Delta S=CZ$, revalue the portfolio to get $L$, and record

$e^{-\theta Q+\psi(\theta)}\mathbf{1}_{\{L>x\}}$.

Average over paths. A practical twisting parameter solves
$\psi'(\theta_x)=x$.

Since $\psi'(\theta)=\mathbb{E}_\theta[Q]$, this samples under a distribution with $\mathbb{E}_{\theta_x}[Q]=x$, moving the rare threshold near the center of the new distribution.

\subsubsection*{Stratified IS}

Stratify the IS estimator by $Q$. For strata $A_1,\ldots,A_K$,

$\mathbb{E}_\theta[e^{-\theta Q+\psi(\theta)}\mathbf{1}_{\{L>x\}}]=\sum_{k=1}^KP_\theta(Q\in A_k)\mathbb{E}_\theta[e^{-\theta Q+\psi(\theta)}\mathbf{1}_{\{L>x\}}\mid Q\in A_k]$.

Under $P_\theta$,
$\psi_\theta(u)=\psi(\theta+u)-\psi(\theta),\qquad \hat\phi_\theta(u)=e^{\psi(\theta+iu)-\psi(\theta)}$.

Use inversion to find $F_\theta(x)=P_\theta(Q\le x)$ and stratum endpoints $a_k$ with desired probabilities, $A_k=(a_{k-1},a_k)$.

Conditional sampling uses acceptance-rejection: generate $Z\sim N(\mu(\theta),\Sigma(\theta))$, compute $Q$, assign to the stratum containing $Q$, keep it if that stratum still needs samples, otherwise discard. This can be worthwhile because generating $Z$ is cheap relative to revaluation.

With samples $(Q_{kj},L_{kj})$ in stratum $k$, estimate

$\sum_{k=1}^KP_\theta(Q\in A_k)\frac1{n_k}\sum_{j=1}^{n_k}e^{-\theta Q_{kj}+\psi(\theta)}\mathbf{1}_{\{L_{kj}>x\}}$.

\subsection{Numerical Example Notes}

The lecture tests short option portfolios on 10 or 100 GBM assets. Prices, deltas, and gammas use Black-Scholes; volatility is 30\%, $r=5\%$, and most cases use uncorrelated assets. Thresholds are set by

$x=\mathbb{E}[Q]+x_{\mathrm{std}}\sqrt{\operatorname{var}(Q)}=a+\sum_{j=1}^m\lambda_j+x_{\mathrm{std}}\sqrt{\sum_{j=1}^m(b_j^2+2\lambda_j^2)}$.

Reported variance-reduction factors compare ordinary MC with control variates, importance sampling, and importance sampling plus stratification.
